\newpage

% ######################################################################
% PART XI: DEMO SCRIPTS & CASE STUDIES
% ######################################################################
\part{Demo Scripts, Mock Interview \& Case Studies}

% ######################################################################
% SECTION 47: DEMO SCRIPTS
% ######################################################################
\section{Demo Scripts --- 30 Sec, 2 Min, 10 Min}

\subsection{30-Second Version (Elevator Pitch)}

\begin{interviewbox}
\textbf{Script (word-by-word):}

``Maine ek \textbf{multi-agent AI research system} banaya hai. User koi bhi topic deta hai, system \textbf{web search} karta hai (Tavily), data ko \textbf{analyze} karta hai, \textbf{facts verify} karta hai ek dedicated fact-checker agent se, aur \textbf{verified report} generate karta hai --- with revision loop agar claims weak milein. Teen interfaces: \textbf{Streamlit} UI, \textbf{MCP server} (AI assistants ke liye), \textbf{Django REST API}. RAG memory: PDFs + past research \textbf{ChromaDB} mein. Aur ek \textbf{benchmark suite} hai jo single vs multi-agent quality measure karta hai.\textbf{''}
\end{interviewbox}

\subsection{2-Minute Version (With Architecture)}

\begin{interviewbox}
\textbf{Script:}

``Pehle 30-sec wala part, phir:''

\textbf{Architecture:} LangGraph StateGraph --- shared state (TypedDict) jisme har node apna output add karta hai. 4 nodes: \textbf{researcher} (LLM se 3 diverse search queries, Tavily se 15 sources, RAG retrieval), \textbf{analyst} (top-10 sources se themes/contradictions/gaps ka dense brief), \textbf{fact-checker} (closed-world verification --- claims sirf sources se VERIFIED; fail par revision loop max 3 baar), \textbf{writer} (brief + top-8 sources se final Markdown report, ChromaDB memory store).

\textbf{Interfaces:} Ek compiled graph, teen adapters --- Streamlit real-time streaming (\texttt{graph.stream}), MCP tools (Claude Desktop connect), Django REST (programmatic + benchmark).

\textbf{Engineering:} 16 pytest tests with full mocking (\textasciitilde9s), config 3-level secret fallback, graceful degradation har external call par, aur benchmark (depth + verifiability, LLM judge + heuristic fallback). Real runs verified: end-to-end 176s, 15 sources, first-pass fact-check pass.\textbf{''}
\end{interviewbox}

\subsection{10-Minute Version (Live Demo Walkthrough)}

\begin{interviewbox}
\textbf{Structure:}
\begin{enumerate}[label=\arabic*.]
\item 0-1 min: Elevator pitch + ek line problem statement.
\item 1-3 min: Architecture draw karo (whiteboard): 4 layers + data flow.
\item 3-5 min: \textbf{Live demo} --- Streamlit: topic run karo, agent statuses dikho, report + download.
\item 5-6 min: PDF upload + RAG retrieval dikhao (``ab ye knowledge base use hoga'').
\item 6-7 min: \textbf{Benchmark Lab} --- run benchmark, verdict dikhao, history charts.
\item 7-8 min: Django API --- \texttt{curl /api/health/}, register, execute (ya Postman screenshot).
\item 8-9 min: Tests --- \texttt{pytest} chalao (9s), output dikhao.
\item 9-10 min: Code tour --- graph/workflow.py ke 40 lines, ek agent file.
\end{enumerate}

\textbf{Rule:} Kabhi bhi 30 sec se zyada ek cheez par mat ruko --- interviewer ko breadth + depth dono dikho.
\end{interviewbox}

% ######################################################################
% SECTION 48: MOCK INTERVIEW DIALOGUE
% ######################################################################
\section{Mock Interview Dialogue --- Poora Session}

\begin{interviewbox}
\textbf{Interviewer:} ``So, tell me about your project.''

\textbf{You:} (30-sec pitch bol do) ``Maine multi-agent research orchestration system banaya hai...'' (section 47 ka script)

\textbf{Interviewer:} ``Why multi-agent? Couldn't one LLM call do the same?''

\textbf{You:} ``Single call se \textbf{verification} nahi hoti. Research quality 3 cheezon par depend karti hai: \textbf{coverage} (sahi sources mile), \textbf{synthesis} (data se insight bani), aur \textbf{verification} (claims true hain). Ek call mein teeno ek saath nahi hote --- agents specialize karte hain. Aur maine isko \textbf{benchmark} se prove karne ki koshish ki --- depth aur verifiability measure karke. Result interesting tha: ek run mein single-agent jeet bhi gaya --- isliye evaluation zaroori hai, assumptions nahi.''

\textbf{Interviewer:} ``What's the most interesting technical challenge?''

\textbf{You:} (Bug 1 ya Bug 5 ki story --- section 39) ``FastMCP version mismatch se MCP server crash --- root cause nikala, fix kiya. Aur benchmark mein multi-agent haar gayi --- root cause: writer-analyst handoff. Ye debugging + root-cause analysis wali mindset dikhati hai.''

\textbf{Interviewer:} ``How would you scale this?''

\textbf{You:} (Section 25 scaling story) ``Synchronous 176s pipeline ko async karna (task queue), SQLite to PostgreSQL, local ChromaDB to managed vector DB, Streamlit multi-instance, rate limiting + caching. Pehla step hamesha bottleneck measure karna.''

\textbf{Interviewer:} ``Do you have questions for us?''

\textbf{You:} (Section 40 Q64 ke 3 questions)
\end{interviewbox}

% ######################################################################
% SECTION 49: CASE STUDIES
% ######################################################################
\section{Case Studies --- 5 Topics Ka Deep Analysis}

\subsection{Case Study 1: Finance Research}

\begin{conceptbox}
\textbf{Topic:} ``Indian EV market 2026'' \textbf{Pipeline behaviour:}
\begin{itemize}
\item Researcher: sub-queries --- market size, policy incentives, players, supply chain.
\item Analyst: themes (growth drivers, policy, competition), contradictions (market size estimates alag sources mein), gaps (battery import data).
\item Fact-checker: market numbers sources se verify.
\item Writer: Executive summary + investment outlook + policy table.
\end{itemize}
\textbf{Interview value:} Finance domain mein \textbf{verifiability} critical hai --- citations + numbers cross-check. Is project ki fact-checker yahin shine karta hai.
\end{conceptbox}

\subsection{Case Study 2: Academic Research}

\begin{conceptbox}
\textbf{Topic:} ``Attention mechanisms survey'' \textbf{Pipeline behaviour:}
\begin{itemize}
\item RAG: uploaded papers (PDFs) se context --- research\_docs collection.
\item Analyst: 3-5 themes (architectures, benchmarks, limitations).
\item Fact-checker: claims papers se verify (closed-world = sirf uploaded papers).
\item Writer: survey-style report with citations.
\end{itemize}
\textbf{Interview value:} Ye dikhata hai ki RAG (PDFs) + agents ka combination academic literature review mein kaam karta hai --- thesis/project work se connect karo.
\end{conceptbox}

\subsection{Case Study 3: Tech Comparison}

\begin{conceptbox}
\textbf{Topic:} ``Rust vs Go for backend 2026'' \textbf{Pipeline behaviour:}
\begin{itemize}
\item Contradictions strong honge (community opinions alag).
\item Fact-checker: performance claims benchmarks se verify.
\item Writer: comparison tables + use-case recommendations.
\end{itemize}
\textbf{Interview value:} Contradiction handling ka perfect example --- system \textbf{disagreement ko hide nahi karta}, report mein highlight karta hai (``Consensus \& Contradictions'' section). Ye honest AI hai.
\end{conceptbox}

\subsection{Case Study 4: Health Research}

\begin{conceptbox}
\textbf{Topic:} ``AI in medical diagnosis 2026'' \textbf{Pipeline behaviour:}
\begin{itemize}
\item High stakes --- hallucination dangerous. Fact-checker critical: har claim source-backed.
\item Gaps flag: ``clinical trial data missing'' --- writer honest reporting.
\end{itemize}
\textbf{Interview value:} Ye case study batao jab interviewer puchhe ``is system ke limitations kya hain?'' --- medical domain mein human-in-the-loop + regulatory compliance zaroori hota. Aap limitations jaante ho --- ye maturity dikhata hai.
\end{conceptbox}

\subsection{Case Study 5: How to Answer ``Improve This System''}

\begin{interviewbox}
\textbf{Question:} ``Is system ko production-ready banane ke liye kya karenge?''

\textbf{Answer (prioritized roadmap):}
\begin{enumerate}[label=\arabic*.]
\item \textbf{Async pipeline} --- Celery/Redis, results stored, polling/webhook (176s sync request khatam).
\item \textbf{Evaluation infra} --- 20-topic benchmark suite + CI mein regression check.
\item \textbf{Structured outputs} --- Pydantic schemas har agent ke (free text $\rightarrow$ typed).
\item \textbf{Observability} --- LangSmith tracing, per-node latency, cost dashboards.
\item \textbf{Security} --- rate limiting, input sanitization, CORS allowlist, HTTPS.
\item \textbf{Data layer} --- PostgreSQL + managed vector DB + backups.
\item \textbf{Multi-model fallback} --- Groq down $\rightarrow$ OpenAI.
\item \textbf{Caching} --- same-topic results + search result dedup.
\end{enumerate}
\textbf{Ordering logic:} Pehle reliability (async), phir quality measurement, phir scale --- har point justify karo.
\end{interviewbox}

\begin{memorybox}
\textbf{Document ka closing thought:}

Ye 150-page document aapke saath interview tak rahega. \textbf{Jo sabse important hai:} concepts itne samjho ki bina notes ke bolo --- ye document notes nahi, \textbf{understanding ka tool} hai. Har section khud bol ke practice karo. Numbers yaad rakho (16 tests, 4 agents, 3 interfaces, 176s, 15 sources). Stories ready rakho (bugs, benchmark loss, fixes). Aur sabse badi baat --- \textbf{genuinely interested} raho: ye project tumhara hai, tumne banaya hai, aur isse pyar karte ho. Wo energy interview mein dikhti hai.

\textbf{All the best, bhai!} \faGraduationCap \faThumbsUp
\end{memorybox}
